\chapter{Calendar-year claims development results}

Merz and W\"uthrich condition on the cells available after a calendar diagonal
\cite{MerzWuethrich2008}. This chapter gives that filtration directly and keeps
their Appendix A.1 linear approximation separate from the exact probability
statements.

\begin{definition}[Calendar-year information]\label{def:calendarSigma}\lean{VerifiedReserving.RandomTriangle.calendarSigma, VerifiedReserving.RandomTriangle.otherCalendar}\leanok
At calendar time $t$,
\[
 \mathcal D^{\mathrm{cal}}_t
 = \sigma\{C_{i,j}: i<n,\ i+j\le t\}.
\]
The companion sigma-algebra \texttt{otherCalendar} contains the observed cells
from every row except a specified accident year.
\end{definition}

\begin{theorem}[The calendar triangle splits at one row]\label{thm:calendarSplit}\lean{VerifiedReserving.RandomTriangle.calendarSigma_mono, VerifiedReserving.RandomTriangle.calendarSigma_eq_sup, VerifiedReserving.condExp_calendarSigma_eq_rowSigma}\leanok
\uses{def:calendarSigma, thm:independenceWitnessRows}
At time $i+k$, the calendar triangle is row $i$ through development $k$ joined
with the observed parts of the other rows. Under accident-year independence,
conditioning a row-$i$ random variable on this calendar triangle is the same as
conditioning it on row $i$ through $k$.
\end{theorem}
\begin{proof}\leanok Split the generating joins at row $i$, then apply conditioning on an independent enlargement.\end{proof}

\begin{definition}[True calendar-year CDR]\label{def:calendarCDR}\lean{VerifiedReserving.RandomTriangle.calendarTrueCDR}\leanok
The true claims development result over $t\to t+1$ is
\[
 E[C_{i,n-1}\mid\mathcal D^{\mathrm{cal}}_t]
 -E[C_{i,n-1}\mid\mathcal D^{\mathrm{cal}}_{t+1}].
\]
\end{definition}

\begin{theorem}[Exact calendar-year CDR]\label{thm:calendarCDR}\lean{VerifiedReserving.condExp_calendarTrueCDR_eq_zero, VerifiedReserving.calendarTrueCDR_eq_of_rows, VerifiedReserving.condExp_sq_calendarTrueCDR_of_rows, VerifiedReserving.condVar_calendarTrueCDR_of_rows}\leanok
\uses{def:calendarCDR, thm:calendarSplit, def:Mack1, def:Mack3}
The true calendar CDR has conditional mean zero. Under the row-conditioned
Mack assumptions and independent accident years, at time $i+k$ it equals
\[
 \left(\prod_{j=k+1}^{n-2} f_j\right)
 (f_kC_{i,k}-C_{i,k+1}),
\]
and its conditional second moment and conditional variance are
$\left(\prod_{j=k+1}^{n-2}f_j\right)^2\sigma_k^2C_{i,k}$.
\end{theorem}
\begin{proof}\leanok The tower property gives the mean-zero result. The split theorem reduces the other two statements to the row filtration and (M1row), (M3row).\end{proof}

\begin{definition}[Appendix A.1 boundary]\label{def:mwA1}\lean{VerifiedReserving.mwA1ProductIncrement, VerifiedReserving.mwA1LinearApprox, VerifiedReserving.mwA1Remainder}\leanok
For finite increments $a_j$, the source replaces
$\prod_j(1+a_j)-1$ by the first-order expression $\sum_j a_j$.
Lean records the exact product, the linear approximation, and their difference
as three separate definitions. No equality or statistical error bound is built
into the approximation.
\end{definition}

\begin{theorem}[Exact boundary identities]\label{thm:mwA1}\lean{VerifiedReserving.mwA1ProductIncrement_eq_linear_add_remainder, VerifiedReserving.mwA1LinearApprox_le_productIncrement, VerifiedReserving.mwA1Remainder_pair}\leanok
\uses{def:mwA1}
The product increment is the linear approximation plus its exact remainder.
For nonnegative increments the linear term is a lower bound, as stated after
(A.1), and for two increments the discarded term is exactly $ab$.
\end{theorem}
\begin{proof}\leanok Ring algebra and the finite Weierstrass product inequality.\end{proof}

\begin{theorem}[Finite CDR and MSEP witness]\label{thm:cdrWitness}\lean{VerifiedReserving.IndependenceWitness.mack3Row, VerifiedReserving.IndependenceWitness.mack3_from_rows, VerifiedReserving.IndependenceWitness.condMsep_witness, VerifiedReserving.IndependenceWitness.condMsepTotal_witness, VerifiedReserving.IndependenceWitness.calendarTrueCDR_three_nondegenerate}\leanok
\uses{thm:calendarCDR, thm:msepRows, thm:msepTotalRows, thm:independenceWitnessRows}
The eight-outcome independent-row model satisfies (M3row), instantiates the
exact single-row and total MSEP wrappers, and has nonzero calendar CDR on three
distinct updates: rows $0$, $1$, and $2$ reveal their independent first-step
shocks at calendar times $0$, $1$, and $2$.
\end{theorem}
\begin{proof}\leanok Each row shock is Rademacher. The first residual has the prescribed second moment and every later residual vanishes.\end{proof}
